\section{攻击面分层：模型、系统、环境与工具链}

上一章已经把安全目标从“保护对象”扩展到“约束行为”，本章进一步回答这些目标会在哪些边界被突破。阅读四层攻击面时，不应把它们理解成互斥标签：真实攻击通常从环境中的不可信内容开始，穿过模型的指令解释，利用系统缺少参数校验，最后借工具权限产生后果。分层的价值是让每一跳都有 owner、检测信号和确定性阻断点。

课程将攻击面拆解为四层：模型层漏洞、系统层漏洞、环境层可利用点、工具链供应链风险。该分层有两个作用：一是避免把所有问题都归因于 prompt injection；二是指导防御资源优先级分配。

\begin{knowledgebox}{模型层典型风险}
包括 prompt engineering attacks、jailbreak、数据投毒导致的错误遵循、拒绝服务型输入（超长上下文、资源消耗诱导）等。模型层风险是基础，但不是全部。
\end{knowledgebox}

\begin{knowledgebox}{系统层典型风险}
包括会话状态机设计缺陷、审批逻辑缺失、工具参数校验不足、跨组件信任链错误、日志与审计盲区。
\end{knowledgebox}

\begin{knowledgebox}{环境层与工具链风险}
包括恶意网页、恶意邮件、恶意 issue、第三方 API 污染、工具插件供应链被劫持等。攻击者往往不直接攻击模型，而是污染其感知来源与执行目标。
\end{knowledgebox}

\begin{warningbox}{为什么 ``只做模型加固'' 不够}
即使模型本身更强健，只要执行层缺少权限隔离，攻击者仍可通过低成本注入驱动高危工具调用，造成真实损害。
\end{warningbox}

例如，一个网页上的间接指令本身属于环境层输入；模型错误服从属于模型层失守；host 把自由文本直接映射为 shell argument 属于系统层缺陷；执行插件使用长期管理员凭证则属于工具链与权限设计问题。只修其中任何一处都可能降低风险，但根因分析必须保存完整链条，否则团队会把下一次变体当成“全新漏洞”重复救火。更好的安全 backlog 不是按攻击名称排队，而是按可复用断点排序：typed interface、parameterization、capability token、transaction、sandbox 和 provenance 往往能同时消除多个攻击家族。

\begin{knowledgebox}{跨层审计的三个问题}
\begin{enumerate}
    \item 哪个低信任输入首次影响了高信任决策？
    \item 哪个组件把建议升级成了有副作用的动作？
    \item 若前两层都失守，哪一层本应限制权限或支持回滚？
\end{enumerate}
回答这三问，可以把“模型被骗了”改写成可修复的系统缺口。
\end{knowledgebox}

\subsection{本章小结}
Agent 风险不是单层风险，而是跨层攻击链。安全工作必须从 ``修一个漏洞'' 转向 ``切断可行攻击路径''。

